\section{Conclusion}

This work demonstrates that interpretability-guided interventions can effectively mitigate bias in large language models by targeting internal representations. Using Sparse Autoencoders to decompose Phi-3-mini-4k-instruct's hidden states, we identified bias-sensitive features and applied encoder-direction steering, achieving a $10.3\%$ relative improvement in anti-stereotype accuracy on CrowS-Pairs ($36.8\% \to 40.58\%$) while maintaining performance on StereoSet. Our primary finding is that \emph{encoder directions} ($\mathbf{E}_k$)—the directions that cause feature activations—significantly outperform both decoder directions ($\mathbf{D}_{:k}$) and direct activation steering for bias mitigation. This suggests that encoder-based interventions provide better control for any interpretability-guided steering task, as they remove bias-aligned components at their source rather than counteracting reconstructions.

Our approach offers several advantages over outer alignment techniques such as DPO and RLHF. While these methods optimize observable behavior through reward signals or preference comparisons, they leave internal representational mechanisms unchanged, allowing biases to persist in hidden layers and propagate through the model. In contrast, our interpretability-guided method directly identifies and modifies bias-aligned components in activation space, providing transparency into \emph{where} and \emph{how} bias manifests. Moreover, our intervention operates at inference time without requiring model retraining, human feedback loops, or preference datasets, making it more practical for deployment scenarios where retraining is costly or infeasible. By targeting the root cause of bias in representations rather than merely suppressing biased outputs, we achieve more fundamental and interpretable bias mitigation.

The superior performance of encoder-direction steering can be understood through the geometric structure of SAE feature spaces. Encoder rows define directions in activation space that directly correspond to feature firing, making them the appropriate handle for targeted removal. In contrast, decoder columns are optimized for reconstruction, which may not align with the causal pathways through which bias influences outputs. Our method incurs a modest utility cost (WikiText-2 perplexity increases by $11.5\%$), but this tradeoff is favorable given the bias reduction achieved. The effectiveness varies across benchmarks: CrowS-Pairs, with localized stereotypical cues, shows significant improvement, while StereoSet, requiring broader contextual reasoning, remains relatively unchanged.

Several limitations constrain our findings: experiments are limited to a single model and layer, we lack category-level error analysis, and our comparison to training-time debiasing methods is limited. Future work should explore utility-preserving constraints, multi-layer interventions, robustness analysis across model families, and deeper comparisons to dataset-level baselines. More broadly, our work contributes to moving beyond alignment-by-behavior toward alignment-by-representation, demonstrating that interpretable feature discovery can inform actionable interventions for bias mitigation without model retraining.
